\subsection{Event Instruction Tuning for Composable Events}

Recent interactive world models expose event-related controls at different
granularities.
LingBot-World~\citep{gao2026lingbotworld} demonstrates promptable global and
local world events, HY-WorldPlay 1.5~\citep{hyworld2025,sun2025worldplay}
supports text-triggered dynamic events during streaming generation, and
Yume-1.5~\citep{mao2025yume15} introduces text-controlled event editing.
Matrix-Game 3.0~\citep{wang2026matrixgame} focuses instead on action-conditioned
long-horizon memory and real-time streaming.
Related spatially controllable video-generation work, such as
Omni-Effects~\citep{mao2025omnieffects}, studies compositional visual effects
specified by category and location, but does not target persistent interactive
world events under navigation.
However, existing public systems do not provide explicit \emph{composable
events}---multiple objects appearing, acting, and interacting within the same
generation under structured event instructions
(Table~\ref{tab:event-comparison}).
Composable events are essential for realistic world simulation, where
meaningful state changes rarely involve a single object in isolation:
a traffic scene requires pedestrians, vehicles, and signals to respond
concurrently, and an indoor scene may involve multiple characters acting and
reacting to one another.
\method{} introduces an Event Instruction Tuning stage that closes this gap:
users describe the coarse region or relation in which each object appears,
what it does, and how objects interact, and the model responds to all specified
events in a single forward pass.

\begin{table}[!htbp]
\caption{Comparison of event-control capabilities across interactive world
models.
Composable events subsume multi-object and multi-event settings: the model
handles multiple objects with distinct actions and mutual interactions in one
generation.}
\label{tab:event-comparison}
\centering
\renewcommand{\arraystretch}{1.2}
\scriptsize
\setlength{\tabcolsep}{2pt}
\begin{tabularx}{\linewidth}{@{}p{0.95in}*{5}{>{\centering\arraybackslash}X}@{}}
\toprule
\rowcolor{amapblue}
\textcolor{white}{\textbf{Model}} &
\textcolor{white}{\textbf{\shortstack{Promptable\\Events}}} &
\textcolor{white}{\textbf{\shortstack{Object-Level\\Events}}} &
\textcolor{white}{\textbf{\shortstack{Region-Guided\\Events}}} &
\textcolor{white}{\textbf{\shortstack{Multi-Entity\\Composition}}} &
\textcolor{white}{\textbf{\shortstack{Inter-Object\\Interaction}}} \\
\midrule
\rowcolor{TableRowAlt}
LingBot-World & \cmark & \cmark & \pmark & \pmark & \xmark \\
HY-WorldPlay 1.5 & \cmark & \cmark & \xmark & \pmark & \xmark \\
\rowcolor{TableRowAlt}
Matrix-Game 3.0 & \xmark & \xmark & \xmark & \xmark & \xmark \\
Yume-1.5 & \cmark & \cmark & \xmark & \xmark & \xmark \\
\rowcolor{TableRowAlt}
\textbf{\method{}} & \cmark & \cmark & \cmark & \cmark & \cmark \\
\bottomrule
\end{tabularx}
\vspace{2pt}
\scriptsize{\pmark indicates qualitative or partial support without structured
per-object event instructions. Object-level events denote single object/category
events; multi-entity composition requires distinct entities and actions in one
structured instruction. Region-guided events refer to coarse regions or relative
placement rather than coordinate-level localization.}
\end{table}

\paragraph{\bf{Training.}}
Using the event instruction data described in Section~\ref{sec:event-data},
we fine-tune the full DiT while keeping the architecture unchanged.
Event semantics enter exclusively through the text-conditioning interface,
where structured event instructions are rendered as natural-language prompts
covering the global scene and per-entity dynamics.
The tuning mixture combines event-instruction samples with non-event training
clips, which preserves the model's general world-generation capability while
adding responsiveness to promptable events.
We use conservative updates and strict gradient clipping to avoid disrupting
the pretrained visual and motion priors.
